% -------------------------------------------------------------
\section{Conclusion and Outlook}
% -------------------------------------------------------------

The Compton effect marks the turning point between classical physics and quantum field theory.
For the first time, it became evident that light is not only a wave
but also a carrier of momentum —
and that the conservation of energy and momentum can only be explained
through the concept of the photon.

However, the Compton formula remains a model that still treats
the photon and the electron as particles in a classical collision.
In reality, it describes a quantized interaction
between two fields — the electromagnetic field and the electron field.

\begin{DidacticBox}[From Formula to Field]
	Quantum electrodynamics no longer describes the Compton effect
	in terms of individual photon trajectories,
	but as an interaction between fields.
	Compton scattering arises there as the exchange of \emph{virtual photons}.
\end{DidacticBox}

The Compton effect unites experimental observation,
relativity, and quantum mechanics within a single equation.
It thus forms the bridge to quantum electrodynamics —
the comprehensive theory of light and its interaction with matter.
